% =====================================================================
%  24.1  Recursion at the Balance -- building the centroid tree.
%  Exports: sec:c18-tree  def:c18-decomp  thm:c18-height
%           fig:c18-centtree  lst:c18-build  prob:c18-321c  lst:c18-321c
%  Solved CF 321C "Ciel the Commander" (2100), live-verified via the
%  problemset API (tags constructive algorithms / dfs / divide and conquer
%  / greedy / trees).  Official samples read live:
%    S1: n=4, star centred at 1 (edges 1-2,1-3,1-4) -> "A B B B".
%    S2: n=10, the path 1-2-...-10 -> "D C B A D C B D C D" (SPECIAL JUDGE:
%        any valid rank assignment is accepted, so the centroid-depth answer
%        need not equal the printed sample).
%  Hand-verified here: S1 centroid = 1 (remove 1 -> three singleton pieces),
%  cd(1)=0 -> 'A', cd(2)=cd(3)=cd(4)=1 -> 'B' => "A B B B", matching S1.
%  Back-cites Ch.12 thm:c09-centfind / lst:c09-cent for the find; this
%  section NEVER re-derives existence -- it begins at "delete it and recurse".
% =====================================================================
\section{Recursion at the Balance}\label{sec:c18-tree}

\begin{previouslybox}[Ch.~12 handed you one centroid and stopped on purpose]
Chapter~12 proved that every tree has a node \(c\) whose removal leaves no piece
bigger than \(\lfloor n/2\rfloor\), and that one DFS finds it
(\(\thmref{c09-centfind}\), \(\coderef{c09-cent}\)). It then stopped: \emph{one}
centroid, \emph{one} cut. This chapter opens the door it left ajar. A centroid is
not a fact to learn once -- it is the pivot of a divide-and-conquer. Delete \(c\),
then do the \emph{same thing} inside every leftover piece. Because each piece has
at most \(n/2\) nodes, the recursion is shallow; the tree of "who split me from
whom" it leaves behind -- the \term{centroid tree} (indexed once in the opener) --
is the object every later section stands on.
\end{previouslybox}

\begin{topicintro}
The whole decomposition is three lines of intent. Given a live component with a
chosen entry node: (1) find its centroid \(c\) with Chapter~12's find; (2)
\emph{mark \(c\) deleted} so no deeper call can step on it; (3) recurse
independently on each neighbour of \(c\) that is still live -- each such neighbour
roots one surviving piece. The centroid found at a deeper call declares the
centroid one level up as its \term{centroid tree} parent, writing
\(\mathit{par}_c\). The recursion's shape is the payoff: removing a centroid caps
every piece at \(|S|/2\), so a component's size at least halves every level --
the centroid tree has height \(\le\lceil\log_2 n\rceil\) (\(\thmref{c18-height}\))
and the build is \(\bigO{n\log n}\).

\textbf{Prerequisites.} Chapter~12 (\chref{c09}) -- the centroid definition
(\(\defref{c09-centroid}\)), existence (\(\thmref{c09-centexist}\)) and the
\(\bigO{n}\) find (\(\thmref{c09-centfind}\)), all used here as a black box. The
subtree-size DFS and the "upward piece \(n-\size(u)\)" discipline are also
Chapter~12's; nothing about distances appears in this section.

\textbf{After this section you can}
\begin{itemize}[nosep]
  \item run the cut-and-recurse loop on a small tree by hand and read off
        \(\mathit{par}_c\) and the centroid depth \(\mathit{cd}\); \hfill\lv{3}
  \item write the recursive build with a correct deleted-marker discipline -- a
        removed centroid is never re-entered -- in \(\bigO{n\log n}\); \hfill\lv{4}
  \item prove the height bound \(\mathit{cd}\le\lceil\log_2 n\rceil\) from the
        "every level halves" fact; \hfill\lv{4}
  \item solve CF~321C: the centroid depth \emph{is} the rank, so the build prints
        the answer directly. \hfill\lv{4}
\end{itemize}
\end{topicintro}

\begin{recognitionbox}[Pattern card -- "a property true along the whole recursion"]
\reckey{statements whose answer is a per-node number you can only get by
recursively re-centring -- \emph{rank every node so that two nodes of the same
rank are watched by a higher one} (CF~321C), \emph{sum over all nodes of their
centroid-tree depth}, or \emph{any quantity defined by "look at the balance point,
solve each side the same way, combine"}. The tell is a global, all-pairs or
whole-tree property with no natural fixed root.}
\recuse{
  \item the decomposition \emph{itself} is the output -- build \(\mathit{par}_c\), read \(\mathit{cd}\) (\(\secref{c18-tree}\));
  \item a path statistic a single rooted DP cannot fold -- claim paths at the centroid (\(\secref{c18-count}\));
  \item an all-component distance index reached from one touched node (\(\secref{c18-index}\)).
}
\recnotuse{
  \item just "print a centroid" -- that is Chapter~12's \(\thmref{c09-centfind}\), no recursion needed;
  \item a range update/query along one fixed path -- HLD (\chref{c16});
  \item a value that is a state per node with no whole-tree/all-pairs dimension -- tree DP (\chref{c10}).
}
One-line test: \emph{the question survives "cut at the balance point and solve both
sides identically" -- then recurse on the centroid, don't root once and hope.}
\end{recognitionbox}

% ---------------------------------------------------------------
\subsection{The decomposition, and the tree it leaves behind}
% ---------------------------------------------------------------

\begin{definition}{Centroid decomposition and the centroid tree}{c18-decomp}
The \term{centroid decomposition} of \(T\) is the recursion
\(\mathit{solve}(S)\): for a live component \(S\), let \(c=\mathit{cen}(S)\)
(\(\defref{c09-centroid}\)); set \(\mathit{par}_c(c)\) to the centroid that
created \(S\); mark \(c\) deleted; recurse on each connected component of
\(S-c\). The \term{centroid tree} is the tree on the same \(n\) nodes whose
parent of \(u\) is \(\mathit{par}_c(u)\); \(\mathit{cd}(u)\) is \(u\)'s depth in
it, and the original tree \(T\) and the centroid tree share their vertex set but
almost never their edges.
\end{definition}

\begin{insightbox}[Why deleting is the whole discipline]
The recursion never rebuilds the tree; it \emph{marks the centroid dead} and lets
the size-DFS treat every dead node as a wall. The moment a centroid is skipped
rather than deleted, a deeper call routes \emph{through} it into a piece it has
already claimed, a component gets processed twice, and the \(\le n/2\) guarantee
-- hence the \(\bigO{\log n}\) depth -- evaporates. The single line
\(\mathit{dead}[c]=\)~\code{true} is the difference between an
\(\bigO{n\log n}\) build and a re-processing \(O(n^2)\) one; the Common Mistakes
station at the chapter end makes you name the symptom.
\end{insightbox}

The running tree \(T\) of \(\figref{c18-tree}\) (opener) decomposes as follows;
the cuts are exactly the top centroid \(4\) and the three pieces it strands.

\begin{tracebox}[Decomposing \(T\) by hand -- component, its centroid, the pieces it leaves]
\textbf{Whole tree}, \(|S|=15\): centroid \(4\) (remove it leaves \(4,4,6\), all
\(\le 7\)); \(\mathit{par}_c(4)=0\), \(\mathit{cd}(4)=0\). Recurse on the pieces
\(A,B,C\).\\
\textbf{\(A=\{1,5,6,11\}\)}, \(|A|=4\): edges \(1\!-\!5,1\!-\!6,6\!-\!11\); centroid
\(1\) (pieces \(1,2\le 2\); ties with \(6\), keep \(1\)); \(\mathit{par}_c(1)=4\),
\(\mathit{cd}=1\). Pieces: \(\{5\}\to\mathit{par}_c(5)=1\); \(\{6,11\}\to\) centroid
\(6\), \(\mathit{par}_c(6)=1\), then \(\{11\}\to\mathit{par}_c(11)=6\).\\
\textbf{\(B=\{2,7,8,12\}\)}, \(|B|=4\): centroid \(2\) (ties with \(8\));
\(\mathit{par}_c(2)=4\). \(\{7\}\to7\); \(\{8,12\}\to\) centroid \(8\), \(\{12\}\to
\mathit{par}_c(12)=8\).\\
\textbf{\(C=\{3,9,10,13,14,15\}\)}, \(|C|=6\): centroid \(3\) (pieces \(2,3\le 3\);
ties with \(10\)); \(\mathit{par}_c(3)=4\). \(\{9,13\}\to\) centroid \(9\),
\(\mathit{par}_c(13)=9\); \(\{10,14,15\}\) is the path \(10\!-\!14\!-\!15\), centroid
\(14\), \(\mathit{par}_c(14)=3\), then \(\{10\},\{15\}\) with parent \(14\).\\
\textbf{Read off:} depths \(0\!:\!\{4\}\); \(1\!:\!\{1,2,3\}\);
\(2\!:\!\{5,6,7,8,9,14\}\); \(3\!:\!\{10,11,12,13,15\}\). Max depth \(3\le
\lceil\log_2 15\rceil=4\).
\end{tracebox}

\begin{figure}[H]
\centering
\begin{tikzpicture}[diagram, scale=0.78, transform shape,
   ct/.style={circle, draw=primaryblue!70, fill=primaryblue!8, line width=0.7pt,
              minimum size=6mm, font=\scriptsize\bfseries, inner sep=1pt},
   root/.style={circle, draw=warningred, fill=warningred!18, line width=1.1pt,
                minimum size=7mm, font=\footnotesize\bfseries, inner sep=1pt}]
  \node[root] (n4) at (0,3) {4};
  \node[ct] (n1) at (-6,2) {1}; \node[ct] (n2) at (0,2) {2}; \node[ct] (n3) at (6,2) {3};
  \node[ct] (n5) at (-7.2,1) {5}; \node[ct] (n6) at (-4.8,1) {6};
  \node[ct] (n7) at (-1.2,1) {7}; \node[ct] (n8) at (1.2,1) {8};
  \node[ct] (n9) at (4.8,1) {9}; \node[ct] (n14) at (7.6,1) {14};
  \node[ct] (n11) at (-4.8,0) {11}; \node[ct] (n12) at (1.2,0) {12};
  \node[ct] (n13) at (4.8,0) {13}; \node[ct] (n10) at (6.6,0) {10}; \node[ct] (n15) at (8.4,0) {15};
  \draw (n4)--(n1) (n4)--(n2) (n4)--(n3);
  \draw (n1)--(n5) (n1)--(n6) (n6)--(n11);
  \draw (n2)--(n7) (n2)--(n8) (n8)--(n12);
  \draw (n3)--(n9) (n3)--(n14) (n9)--(n13) (n14)--(n10) (n14)--(n15);
  \node[annot, font=\scriptsize, text=darkgray, anchor=east] at (-8.6,3) {\(\mathit{cd}=0\)};
  \node[annot, font=\scriptsize, text=darkgray, anchor=east] at (-8.6,1) {\(\mathit{cd}=2\)};
\end{tikzpicture}
\caption{The \term{centroid tree} of \(T\), built by \(\secref{c18-tree}\)'s
recursion. Its edges are \(\mathit{par}_c\) relations, not \(T\)'s edges: the top
centroid \(4\) (red) roots the tree, \(1,2,3\) are its three piece-centroids, and
the remaining nodes fill depths 2 and 3. The height of \(3\) obeys
\(\thmref{c18-height}\); CF~321C prints exactly this tree, reading each rank as
\(\text{'A'}+\mathit{cd}\).}
\label{fig:c18-centtree}
\end{figure}

\begin{lstlisting}[caption={The recursive centroid-tree build over a fixed tree. The centroid finder is Ch.12's walk-to-heavy (thm:c09-centfind) used as a black box; the dead[] marker keeps pieces out of already-claimed nodes. ASCII only}, label={lst:c18-build}]
// Global: adj[], n. Outputs par_c[1..n] (centroid parent, 0 = top) and cd[1..n].
const int MAXN = 200005;
vector<int> adj[MAXN];
bool dead[MAXN];
int sz[MAXN], par_c[MAXN], cd[MAXN];

int dfs_size(int u, int p){                    // sizes inside this LIVE component
    sz[u] = 1;
    for(int v : adj[u])
        if(v != p && !dead[v]) sz[u] += dfs_size(v, u);
    return sz[u];
}
int find_centroid(int u, int p, int tot){      // Ch.12 thm:c09-centfind, walk-to-heavy
    bool moved = true;
    while(moved){
        moved = false;
        for(int v : adj[u])
            if(v != p && !dead[v] && sz[v] > tot/2){ p = u; u = v; moved = true; break; }
    }
    return u;                                  // here every piece is <= tot/2
}
void decompose(int entry, int parent_centroid){
    int tot = dfs_size(entry, -1);             // entry roots one live piece
    int c   = find_centroid(entry, -1, tot);   // Ch.12's find
    par_c[c] = parent_centroid;                // c's parent in the centroid tree
    cd[c]    = (parent_centroid == 0) ? 0 : cd[parent_centroid] + 1;
    dead[c]  = true;                           // DELETED: a wall for every deeper call
    for(int v : adj[c])                        // each live neighbour roots a surviving piece
        if(!dead[v]) decompose(v, c);
}
// Build the whole centroid tree with:  decompose(1, 0);
\end{lstlisting}

\begin{theorem}{The centroid tree has logarithmic height}{c18-height}
For any tree on \(n\ge 1\) nodes, the centroid tree produced by
\(\secref{c18-tree}\) has height \(\le\lceil\log_2 n\rceil\): every node satisfies
\(\mathit{cd}(u)\le\lceil\log_2 n\rceil\).
\end{theorem}
\begin{proofbox}
Follow the chain of components containing a fixed node \(u\): the whole tree has
\(n\) nodes; after the top centroid is deleted, \(u\)'s component has at most
\(\lfloor n/2\rfloor\) nodes because a centroid caps every piece at half
(\(\defref{c09-centroid}\)); each further level re-centres a component and again
caps \(u\)'s piece at half its parent's size. So the component holding \(u\) has
size \(\le n,\ \lfloor n/2\rfloor,\ \lfloor n/4\rfloor,\dots\) reaching \(1\) after
at most \(\lceil\log_2 n\rceil\) halvings -- and each halving corresponds to
exactly one centroid on \(u\)'s centroid-ancestor chain. Hence
\(\mathit{cd}(u)\le\lceil\log_2 n\rceil\). Two-centroid ties (\(\thmref{c09-centexist}\))
do not affect the bound: whichever of the pair is chosen, both cap every piece at
\(\le n/2\), which is all the induction uses. \(\Box\)
\end{proofbox}

\begin{complexitybox}[The build]
\textbf{Time} \(\bigO{n\log n}\): a node lies in exactly one live component per
centroid level, so it is visited by \(\mathit{dfs\_size}\)/\(\mathit{find\_centroid}\)
once per level, and there are \(\le\lceil\log_2 n\rceil\) levels
(\(\thmref{c18-height}\)). \textbf{Space} \(\bigO{n}\) for \(\mathit{par}_c\),
\(\mathit{cd}\), \(\mathit{sz}\), \(\mathit{dead}\). On a chain the naive
"recolor and recurse on a rooted subtree" would be \(\bigO{n^2}\); the
\(\le n/2\) cut is precisely what collapses it to \(\bigO{n\log n}\) -- the same
halving \(\thmref{c09-centexist}\)'s proof already promised.
\end{complexitybox}

% ---------------------------------------------------------------
\subsection{Problem Lab: CF 321C -- rank every node by its centroid depth}
% ---------------------------------------------------------------

\solvedproblem{c18-321c}{Ciel the Commander}{CF 321C}
  {https://codeforces.com/problemset/problem/321/C}{CF 2100 L4}
  {centroid decomposition; the answer is 'A' + centroid depth; special judge}

\begin{problemstatement}[Ciel the Commander (statement and samples verified against the live problem)]
A tree has \(n\) cities (\(2\le n\le 10^5\)). Assign each city an officer rank, a
letter \(\text{'A}\)–\(\text{'Z}\) ('A' topmost). Rule: for any two distinct cities
\(x,y\) of the \emph{same} rank, the simple path between them must contain a city
of \emph{strictly higher} rank. Output the \(n\) ranks (space-separated), or
\texttt{Impossible!} if none exists. \textbf{Special judge: any valid assignment
is accepted}, so your string may differ from the sample.

\textbf{Examples (official).}
\begin{verbatim}
Input:              Output:
4                   A B B B
1 2
1 3
1 4

Input:              Output:
10                  D C B A D C B D C D
1 2 ... 9 10        (path on 10 nodes; one of many valid answers)
\end{verbatim}
\end{problemstatement}

\begin{attemptgate}[Attempt it yourself first -- \lv{4}, give yourself 25 minutes]
Guess that the topmost rank 'A' goes to a centroid, then ask what rank 'B' can
occupy. \peek{which nodes may share a rank without violating the monitoring rule?}
\nudge{two same-rank nodes are fine if their path holds a higher rank; a centroid's
removal is exactly the thing that separates pieces from each other.}
\fullhint{Rank a node by its centroid-tree depth: \(\text{'A'}+\mathit{cd}(u)\)
(\(\coderef{c18-build}\)). If \(x,y\) have the same depth \(d\), their lowest common
ancestor in the \emph{centroid} tree has depth \(<d\), and that LCA centroid lies on
the original-tree path \(x\text{--}y\) -- the higher-rank monitor the rule wants.
For \(n\le 10^5\), \(\lceil\log_2 n\rceil\le 17<26\), so 'A'..'R' always suffice and
\texttt{Impossible!} never fires on valid input. Build the decomposition and print
\(\text{'A'}+\mathit{cd}[i]\).}
\end{attemptgate}

\solstep{Understanding}
The rank is a function you must place on vertices so that equal values are always
"bridged" by a larger one along their path. That is not a distance problem and not
a fixed-root DP; it is a statement about a \emph{recursive separation} of the tree
into halves -- which is the centroid decomposition. The insight is that
\(\mathit{cd}\) (not \(\size\), not depth in \(T\)) is the right coordinate, and
then the problem is \emph{literally} "run \(\secref{c18-tree}\)'s build and print
the depth."

\solstep{Key Observations}
\begin{observation}
Setting \(\mathrm{rank}(u)=\text{'A'}+\mathit{cd}(u)\) is valid. For same-depth
\(x\ne y\), let \(w\) be their centroid-tree LCA; \(w\) is a proper ancestor of
both so \(\mathit{cd}(w)<\mathit{cd}(x)\), i.e.\ higher rank; and \(w\) is the first
centroid whose cut separated \(x\) and \(y\), which forces \(w\) onto the
\(T\)-path between them.
\end{observation}
\begin{observation}
The bound is always satisfiable here: \(\thmref{c18-height}\) gives
\(\max\mathit{cd}\le\lceil\log_2 10^5\rceil=17\le 25\), so no rank exceeds 'Z' and
the \texttt{Impossible!} branch is unreachable for \(n\le 10^5\).
\end{observation}

\solstep{The Approach Ladder}
\begin{center}\renewcommand{\arraystretch}{1.25}\setlength{\tabcolsep}{3pt}
\begin{tabular}{@{}c p{0.26\linewidth} p{0.30\linewidth} l c@{}}
\toprule
\textbf{\#} & \textbf{Approach} & \textbf{Idea} & \textbf{Time} & \textbf{Verdict}\\
\midrule
1 & \lv{2}\ fixed-root depth & rank by depth in \(T\) & \(\bigO{n}\) & \textcolor{warningred}{\ding{55}\ \textbf{WA}} \\
2 & \lv{3}\ greedy leaf-peel & colour leaves 'Z' upward & \(\bigO{n^2}\) & \tle \\
3 & \lv{4}\ \textbf{centroid depth} & rank\(=\text{'A'}+\mathit{cd}\) (\(\coderef{c18-build}\)) & \(\bigO{n\log n}\) & \ac \\
\bottomrule
\end{tabular}\end{center}
Rung~1 is the seductive trap: depth in \(T\) under a bad root can be \(O(n)\) (a
chain), pushing ranks past 'Z', and it puts two same-rank leaves on one side of
every internal node -- no higher-rank monitor between them. Rung~3's centroid depth
is \(O(\log n)\) (\(\thmref{c18-height}\)) and separates by construction.

\solstep{Solution}
Run \(\mathit{decompose}(1,0)\) from \(\coderef{c18-build}\) to fill \(\mathit{cd}\);
print \(\text{'A'}+\mathit{cd}[i]\) for \(i=1..n\).

\solstep{Implementation}
\begin{lstlisting}[caption={CF 321C driver: build the centroid tree once, then print 'A' + centroid depth per node. Reuses decompose() from lst:c18-build verbatim. ASCII only}, label={lst:c18-321c}]
#include <bits/stdc++.h>
using namespace std;
// ... paste adj[], dead[], sz[], par_c[], cd[], dfs_size, find_centroid,
//     and decompose() exactly as in lst:c18-build ...
int main(){
    int n;  scanf("%d", &n);
    for(int i=0;i<n-1;i++){
        int a,b; scanf("%d %d",&a,&b);
        adj[a].push_back(b); adj[b].push_back(a);
    }
    decompose(1, 0);                       // fills cd[1..n]
    for(int i=1;i<=n;i++)
        printf("%c%c", 'A'+cd[i], i==n ? '\n' : ' ');
    // cd[i] <= 17 <= 25, so a rank is always a letter: never "Impossible!" here.
}
\end{lstlisting}

\solstep{Trace}
\begin{tracebox}[Sample 1: the star \(n=4\), edges \(1\!-\!2,1\!-\!3,1\!-\!4\)]
\(\mathit{decompose}(1,0)\): centroid of the whole star is \(1\) (remove it leaves
\(\{2\},\{3\},\{4\}\), each size \(1\le 2\)); \(\mathit{cd}(1)=0\). Recurse on each
leaf component: centroid of \(\{2\}\) is \(2\), \(\mathit{cd}=1\); likewise \(3,4\).
So \(\mathit{cd}=(0,1,1,1)\) for nodes \((1,2,3,4)\) and the printout is
\(\boxed{\texttt{A B B B}}\) -- matching the official output. On \(T\) of
\(\figref{c18-tree}\) the same code yields
\(1\!:\!\text{B},2\!:\!\text{B},3\!:\!\text{B},4\!:\!\text{A},5\!:\!\text{C},6\!:\!\text{C},
7\!:\!\text{C},8\!:\!\text{C},9\!:\!\text{C},10\!:\!\text{D},11\!:\!\text{D},12\!:\!\text{D},
13\!:\!\text{D},14\!:\!\text{C},15\!:\!\text{D}\) -- the depth layers read straight
off \(\figref{c18-centtree}\).
\end{tracebox}

\begin{warningbox}[Reusing sizes across the recursion -- the signature bug]
\(\mathit{sz}[\cdot]\) is meaningful \emph{only} inside the current component,
rooted at the current \(\mathit{entry}\), \emph{after} the ancestors are dead. If you
compute \(\mathit{sz}\) once at the top and reuse it in deeper calls (skipping the
per-call \(\mathit{dfs\_size}\)), the centroid test is wrong: a stale \(\mathit{sz}\)
counts nodes that belong to an already-deleted sibling. Symptom on a chain: the build
either picks a non-centroid (linear depth, ranks overflow 'Z', WA) or loops. Re-run
\(\mathit{dfs\_size}\) at every \(\mathit{decompose}\) call -- that is the whole
\(\bigO{n\log n}\) budget.
\end{warningbox}

\begin{hardeningbox}
\begin{itemize}[nosep,leftmargin=1.4em]
  \item \textbf{Deletion before recursion.} Set \(\mathit{dead}[c]=\)~\code{true}
        \emph{before} iterating \(c\)'s neighbours, and guard every edge walk with
        \(\mathit{!dead}[\cdot]\); otherwise the recursion steps back onto the pivot.
  \item \textbf{Stack budget.} The centroid recursion is \(O(\log n)\) deep, but the
        inner \(\mathit{dfs\_size}\) is as deep as the \emph{component} (up to \(n\) on
        a chain). For \(n=10^5\) use an explicit stack (Ch.~12's \(\coderef{c09-iter}\))
        or raise the recursion limit -- \(\secref{c18-tree}\)'s listings keep the
        recursive form for legibility.
  \item \textbf{Two-centroid ties.} When a component admits a pair
        (\(\thmref{c09-centexist}\)), any choice is legal -- the height proof uses only
        the \(\le n/2\) cap, so \(\mathit{cd}\) stays bounded and CF~321C stays valid.
  \item \textbf{Special judge.} Do not compare your string to the sample literally;
        verify the rule, or trust \(\thmref{c18-height}\) and the LCA-monitor
        observation.
\end{itemize}
\end{hardeningbox}

\begin{recallbox}[24.1 in one breath]
\emph{Find the centroid (Ch.~12, \(\thmref{c09-centfind}\)), delete it, recurse on
each surviving piece, and record \(\mathit{par}_c\) -- that is the centroid tree
(\(\defref{c18-decomp}\), \(\coderef{c18-build}\)). Every level at least halves a
component, so height \(\le\lceil\log_2 n\rceil\) (\(\thmref{c18-height}\)) and the
build is \(\bigO{n\log n}\). CF~321C is this build seen from the output side: rank
\(=\text{'A'}+\mathit{cd}\), because a same-rank pair's centroid-LCA is a higher-rank
monitor on their path.}
\selftest{
  \item Why is \(\mathit{par}_c\) almost never equal to \(T\)'s parent array? \emph{(\(\mathit{par}_c(u)\) is the centroid that split \(u\)'s component, not its tree parent (\(\figref{c18-centtree}\)).)}
  \item State the invariant that bounds the centroid-tree height. \emph{(Every level caps the component at \(\le\) half (\(\thmref{c18-height}\)).)}
  \item In CF~321C, if two 'B' cities exist, what guarantees a higher rank on their path? \emph{Their centroid-tree LCA has depth \(0\) (rank 'A') and lies on their path (Key Observation~1).)}
}
\end{recallbox}
